\documentclass{article}

\usepackage[margin=1in]{geometry}
\usepackage{graphicx, amsmath, amsthm, amssymb, enumitem, lipsum}
\usepackage[]{hyperref}

\setlist[enumerate]{leftmargin=1.25cm}

\newcommand{\N}{\mathbb{N}}
\newcommand{\Z}{\mathbb{Z}}
\newcommand{\Q}{\mathbb{Q}}
\newcommand{\R}{\mathbb{R}}
\newcommand{\C}{\mathbb{C}}

\newcommand{\rotvert}{\rotatebox[origin=c]{90}{$\vert$}}

\newenvironment{note}[2][Note]{\begin{trivlist}
\item[\hskip \labelsep {\bfseries #1}\hskip \labelsep {\bfseries #2.}]}{\end{trivlist}}

\newenvironment{fact}[2][Fact]{\begin{trivlist}
\item[\hskip \labelsep {\bfseries #1}\hskip \labelsep {\bfseries #2.}]}{\end{trivlist}}

\newenvironment{definition}[2][Definition]{\begin{trivlist}
\item[\hskip \labelsep {\bfseries #1}\hskip \labelsep {\bfseries #2.}]}{\end{trivlist}}

\newenvironment{proposition}[2][Proposition]{\begin{trivlist}
\item[\hskip \labelsep {\bfseries #1}\hskip \labelsep {\bfseries #2.}]}{\end{trivlist}}

\newenvironment{principle}[2][Principle]{\begin{trivlist}
\item[\hskip \labelsep {\bfseries #1}\hskip \labelsep {\bfseries #2.}]}{\end{trivlist}}

\newenvironment{lemma}[2][Lemma]{\begin{trivlist}
\item[\hskip \labelsep {\bfseries #1}\hskip \labelsep {\bfseries #2.}]}{\end{trivlist}}

\newenvironment{theorem}[2][Theorem]{\begin{trivlist}
\item[\hskip \labelsep {\bfseries #1}\hskip \labelsep {\bfseries #2.}]}{\end{trivlist}}

\newenvironment{corollary}[2][Corollary]{\begin{trivlist}
\item[\hskip \labelsep {\bfseries #1}\hskip \labelsep {\bfseries #2.}]}{\end{trivlist}}

\newenvironment{envsection}[1]{\begin{trivlist}
\item[\hskip \labelsep {\bfseries #1}]}{\end{trivlist}}


\begin{document}
\setlength{\abovedisplayskip}{4pt}
\setlength{\belowdisplayskip}{4pt}


\begin{center}
    \Large{\textbf{Homework 3}}\\
    \normalsize February 12 - 19
\end{center}
\vspace{10pt}

\subsection*{Reading}

Feel free to just casually skim/read the following---no need to take notes unless you want to. The readings are meant to complement what we've done so far and should be interesting and enjoyable more than anything.

\begin{envsection}{1.}
    Read \href{http://www.incompleteideas.net/IncIdeas/BitterLesson.html}{The Bitter Lesson} by Rich Sutton.
\end{envsection}

\begin{envsection}{2. Probabilistic Machine Learning:}
    Sections 8.4.1 and 8.4.2 on SGD. The rest of 8.4 is pretty advanced but there looks to be some cool stuff in there---feel free to look around and skim anything that looks interesting.
\end{envsection}

\begin{envsection}{3. CS229 Notes:}
    Read Chapter 8, the introduction and 8.1 (no need to read 8.1.1, unless you're interested). This talks a bit more in depth about generalization and the bias-variance tradeoff, which we read a bit about in R+G 1.5.
\end{envsection}



% \subsection*{Problems}

% \begin{envsection}{Problem 1.}
    
% \end{envsection}


\end{document}